%\documentclass{...}

\input{./def/yf-formatting}

\usepackage{amsfonts, amsmath, amssymb, amsthm}
\usepackage{comment}
\usepackage{graphicx}
\usepackage{ifthen}
\usepackage{latexsym}
%\usepackage{times}
\usepackage[normalem]{ulem}

\input{./def/yf-def}

\def\In{\text{IN}}
\def\out{\text{OUT}}
\def\A{\mathcal{A}}
\def\E{\mathcal{E}}
\def\H{\mathcal{H}}
\def\Q{\mathcal{Q}}
\def\T{\mathcal{T}}
\def\V{\mathcal{V}}
\def\W{\mathscr{W}}

\newboolean{solver}\setboolean{solver}{true}
%\newboolean{solver}\setboolean{solver}{false}
\ifthenelse{\boolean{solver}}{\includecomment{sol}}{\excludecomment{sol}}

\def\vgap{\vspace{2mm}}

\begin{document}

\section*{CE7456: Exercise List 6}
Prepared by Yufei Tao \\

% \begin{center}
%     \uline{Classroom Discussion}
% \end{center}

\boxminipg{0.95\linewidth}{
    These exercises are designed to help you practice the techniques taught in class. Attempt them ``at your own risk'' because they are research-oriented and can be challenging. Solutions will not be provided in written form. I will be happy to act as a collaborator and sounding board --- I look forward to discussing your creative approaches, refining your key steps, and walking through your attempted solutions with you.
}

\extraspacing {\bf Problem 1.} Prove: if a problem can be solved in the RAM model in $T(n)$ time, then it can be solved in the EM model in $T(n)$ I/Os.

\extraspacing {\bf Problem 2.} Let $S$ be a set of $n$ integers. Given an interval $q = [x, y]$, a {\em range query} reports all the elements $e \in S$ such that $e \in q$. Explain how to store $S$ in a data structure of $O(n/B)$ blocks that can answer a range query in $O(\log_B n + k / B)$ I/Os, where $k$ is the number of elements reported.

\extraspacing {\bf Problem 3.} In our lecture notes, we defined the EM model by requiring $M \ge 3B$ (recall that $M$ and $B$ are the memory and block sizes, respectively). However, in the standard definition of the model, $M$ can be as small as $2B$. This exercise is designed to show you that our assumption of $M \ge 3B$ does not lose generality for asymptotic analysis.

\vgap

Let $S$ be a set of $n$ elements given in a disk array. Let $\A$ be an algorithm that takes $S$ as the input and performs $T(M, B)$ I/Os when $M \ge 100B$. Use $\A$ algorithm that performs $O(T(M, B/101))$ I/Os on $S$ when $M = 2B$.

\vgap

(Hint: conceptually chop a block into smaller blocks.)

\extraspacing {\bf Problem 4.} Consider the sorting problem under the MPC model discussed in class. You can assume that the input size $n$ is a multiple of the number $p$ of machines. Suppose that we impose a requirement on the output: the smallest $n/p$ elements must reside on machine 1, the next $n/p$ elements must reside on machine 2, and so on. Under the assumption $n \ge p^2 \ln n$, solve the problem with probability at least $1 - 1/n^2$ with load $O(n/p)$.

\extraspacing {\bf Problem 5.} Consider the {\em triangle listing} problem: given an undirected graph $G = (V, E)$ with $m$ edges, find all the 3-cliques in $G$.

\myitems{
    \item Design an MPC algorithm to solve the problem in one-round with load $O(m / p^{1/3})$.

    \item It is known that there is a constant-round MPC algorithm that solves the problem with load $O(m / p^{2/3})$. Design an EM algorithm for solving the problem in $\tO(m^{1.5} / (B \sqrt{M}) )$ I/Os, where $\tO(.)$ hides a factor polylogarithmic in $m$, $M$, and $B$. Your algorithm only needs to ``see'' each triangle  once in memory, rather than writing it to the disk.
}

\extraspacing {\bf Problem 6.}
\end{document}
